% Part: lambda-calculus
% Chapter: syntax
% Section: alpha

\documentclass[../../../include/open-logic-section]{subfiles}

\begin{document}

\olfileid{lam}{syn}{alp}
\olsection{$\alpha$-રૂપાંતરણ}

$\lambd[x][x]$ અને~$\lambd[y][y]$ વચ્ચે શું સંબંધ છે?
બંને ઓળખ-વિધેય દર્શાવે છે. અલબત્ત, વાક્યરચનાની દૃષ્ટિએ તે
જુદાં પદો છે. ફરક માત્ર બદ્ધ ચલના નામમાં છે; એકમાંના બદ્ધ
ચલનું ``નામ બદલીને'' બીજું મળે છે. તેને
\emph{$\alpha$-રૂપાંતરણ} કહે છે.

\begin{defn}[બદ્ધ ચલનું નામ બદલવું, $\aconvone$]
  પદ~$M$માં~$\lambd[x][N]$ની કોઈ ઘટના હોય, $x\neq y$,
  $y \notin
  \FV{N}$ હોય અને~$\Subst{N}{y}{x}$ વ્યાખ્યાયિત હોય તો
  તે ઘટનાને
  \begin{equation*}
    \lambd[y][\Subst{N}{y}{x}]
  \end{equation*}
  વડે બદલીને~$M'$ મળે તેને \emph{બદ્ધ ચલનું નામ બદલવું}
  કહીએ છીએ અને~$M \redone[\alpha] M'$ લખીએ છીએ.
\sourcecorrection{OLLAM-019}{સ્થિર અંગ્રેજી મૂળની પહેલી
વ્યાખ્યામાં~$x\neq y$ની શરત નથી, પરંતુ આગળની બંને સમકક્ષ
વ્યાખ્યાઓમાં તે છે. અહીં નામ ખરેખર બદલાય તે માટે એ શરત
સ્પષ્ટ મૂકી છે.}
\end{defn}

\begin{defn}[સંબંધની રચના-સુસંગતતા]
  પદો પરનો સંબંધ~$R$ નીચેની શરતો સંતોષે તો તેને
  \emph{રચના-સુસંગત} કહીએ છીએ:
  \begin{enumerate}
  \item જો $R N N'$ હોય તો $R \lambd[x][N] \lambd[x][N']$.
  \item જો $R P P'$ હોય તો $R (PQ) (P'Q)$.
  \item જો $R Q Q'$ હોય તો $R (PQ) (PQ')$.
  \end{enumerate}
\end{defn}

આથી વ્યાખ્યા ફરી આ રીતે કહી શકાય:
\begin{defn}[બદ્ધ ચલનું નામ બદલવું, $\aconvone$]
  \emph{બદ્ધ ચલનું નામ બદલવાનો} સંબંધ~($\redone[\alpha]$)
  પદો પરનો સૌથી નાનો રચના-સુસંગત સંબંધ છે જે નીચેની
  શરત સંતોષે છે:
  \begin{align*}
    &\lambd[x][N] \redone[\alpha] \lambd[y][\Subst{N}{y}{x}] && \text{જો
      $x \neq y$, $y \notin \FV{N}$ હોય} \\
    & &&\text{અને $\Subst{N}{y}{x}$ વ્યાખ્યાયિત હોય}
  \end{align*}
\end{defn}

અહીં ``સૌથી નાનો'' એટલે સંબંધમાં રચના-સુસંગતતા અને વધારાની
શરતથી જરૂરી બનેલી જોડીઓ જ હોય, બીજી કોઈ નહીં. તેથી સંબંધને
આગમનથી પણ વ્યાખ્યાયિત કરી શકાય:

\begin{defn}[બદ્ધ ચલનું નામ બદલવું, $\aconvone$] \ollabel{defn:aconvone}
  \emph{બદ્ધ ચલનું નામ બદલવાનો} સંબંધ~($\aconvone$)
  આગમનથી આ રીતે વ્યાખ્યાયિત થાય છે:
  \begin{enumerate}
  \item જો $N \aconvone N'$ હોય તો $\lambd[x][N] \aconvone
    \lambd[x][N']$. \ollabel{defn:aconvone1}
  \item જો $P \aconvone P'$ હોય તો $(PQ) \aconvone (P'Q)$.
    \ollabel{defn:aconvone2}
  \item જો $Q \aconvone Q'$ હોય તો $(PQ) \aconvone (PQ')$.
    \ollabel{defn:aconvone3}
  \item જો $x \neq y$, $y \notin \FV{N}$ અને~$\Subst{N}{y}{x}$
    વ્યાખ્યાયિત હોય તો
    $\lambd[x][N] \redone[\alpha] \lambd[y][\Subst{N}{y}{x}]$.
    \ollabel{defn:aconvone4}
  \end{enumerate}
\end{defn}

આ વ્યાખ્યાઓ સમકક્ષ છે, પરંતુ તેની સાબિતી અભ્યાસપ્રશ્ન તરીકે
છોડીએ છીએ. હવે પછી આગમનાત્મક વ્યાખ્યા વાપરીશું.

\begin{defn}[$\alpha$-રૂપાંતરણ, $\aconv$]
  \emph{$\alpha$-રૂપાંતરણ}~($\aconv$) એ પદો પરનો
  $\aconvone$ને સમાવતા સ્વવાચક અને પરંપરિત સંબંધોમાં
  સૌથી નાનો સંબંધ છે.
\end{defn}

અગાઉની જેમ, ``સૌથી નાનો'' એટલે સંબંધમાં સ્વવાચકતા,
પરંપરિતતા અને~$\aconvone$ને કારણે જરૂરી બનેલી જોડીઓ જ હોય.
તેમાંથી નીચેની સમકક્ષ વ્યાખ્યા મળે છે:
\sourcecorrection{OLLAM-020}{સ્થિર અંગ્રેજી મૂળની આ
સમજૂતીમાં ``સૌથી નાના'' સંબંધ માટે પરંપરિતતા અને
$\aconvone$નું નામ છે, પરંતુ વ્યાખ્યામાં જરૂરી સ્વવાચકતાનો
ઉલ્લેખ છૂટી ગયો છે. અહીં તે ઉમેર્યો છે.}

\begin{defn}[$\alpha$-રૂપાંતરણ, $\aconv$] \ollabel{defn:aconv}
  \emph{$\alpha$-રૂપાંતરણ}~($\aconv$) આગમનથી આ રીતે
  વ્યાખ્યાયિત થાય છે:
  \begin{enumerate}
  \item જો $P \aconv Q$ અને~$Q \aconv R$ હોય તો
    $P \aconv R$. \ollabel{defn:aconv1}
  \item જો $P \aconvone Q$ હોય તો $P \aconv Q$.
    \ollabel{defn:aconv2}
  \item $P \aconv P$. \ollabel{defn:aconv3}
  \end{enumerate}
\end{defn}

\begin{ex}
  $\lambd[x][f x]$નું~$\lambd[y][f
  y]$માં $\alpha$-રૂપાંતરણ થાય છે, અને ઊલટું પણ.
  અનૌપચારિક રીતે કહીએ તો બંને આર્ગ્યુમેન્ટ સ્વીકારીને તેના
  પર~$f$ લાગુ કરવાનું પરિણામ આપે છે; અહીં~$f$ પરિવેશનો
  ચલ છે.

  $\lambd[x][f x]$નું~$\lambd[x][g
    x]$માં $\alpha$-રૂપાંતરણ થતું નથી. પહેલું અને બીજું પદ
  અનુક્રમે પરિવેશના ચલો~$f$ અને~$g$ તરફ નિર્દેશ કરે છે.
  તેથી એ જુદાં વિધેયો છે: જ્યાં~$f$ અને~$g$ જુદાં હોય એવા
  પરિવેશોમાં તેમનું વર્તન જુદું હોય છે.
\end{ex}

\begin{prob}
  નીચેની પદ-જોડીઓ $\alpha$-રૂપાંતરણથી એકબીજામાં બદલી શકાય છે?
  \begin{enumerate}
  \item $\lambd[x][\lambd[y][x]]$ અને $\lambd[y][\lambd[x][y]]$
  \item $\lambd[x][\lambd[y][x]]$ અને $\lambd[c][\lambd[b][a]]$
  \item $\lambd[x][\lambd[y][x]]$ અને $\lambd[c][\lambd[b][a]]$
  \end{enumerate}
\sourcecorrection{OLLAM-021}{સ્થિર અંગ્રેજી મૂળમાં આ
અભ્યાસપ્રશ્નની બીજી અને ત્રીજી જોડી એકસરખી છે. અનુમાનથી
કોઈ નવી જોડી ન ઘડતાં બંને મૂળ મુજબ જાળવી છે.}
\end{prob}

\begin{lem}\ollabel{lem:fv-one}
  $P \aconvone Q$ હોય તો $\FV{P} = \FV{Q}$.
\end{lem}

\begin{proof}
  $P \aconvone Q$ની !!{derivation} પર આગમનથી.
  \begin{enumerate}
  \item છેલ્લો નિયમ~\olref{defn:aconvone4} હોય તો~$P$
    એ~$\lambd[x][N]$ અને~$Q$ એ
    $\lambd[y][\Subst{N}{y}{x}]$ સ્વરૂપના છે; અહીં
    $x \neq y$, $y \notin \FV{N}$ અને~$\Subst{N}{y}{x}$
    વ્યાખ્યાયિત છે. હવે~$x \in \FV{N}$ છે કે નહીં તે
    મુજબ કિસ્સા કરીએ:
    \begin{enumerate}
    \item $x \in \FV{N}$ હોય તો:
      \begin{align*}
        \FV{\lambd[y][\Subst{N}{y}{x}]} &=
        \FV{\Subst{N}{y}{x}} \setminus \{y\} \\
        & = ((\FV{N} \setminus \{x\}) \cup \{y\}) \setminus \{y\}
         && \text{\olref[sub]{thm:infv} મુજબ} \\
        & = \FV{N} \setminus \{x\} \\
        & = \FV{\lambd[x][N]}
      \end{align*}
    \item $x \notin \FV{N}$ હોય તો:
      \begin{align*}
        \FV{\lambd[y][\Subst{N}{y}{x}]}
        & = \FV{\Subst{N}{y}{x}} \setminus \{y\} \\
        & = \FV{N} \setminus \{x\}
         && \text{\olref[sub]{thm:notinfv} મુજબ} \\
        & = \FV{\lambd[x][N]}.
      \end{align*}
    \end{enumerate}
  \item બીજા ત્રણ કિસ્સા અભ્યાસપ્રશ્ન તરીકે છોડીએ છીએ.
  \end{enumerate}
\sourcecorrection{OLLAM-022}{સ્થિર અંગ્રેજી મૂળની આ
સાબિતીમાં~$x\in FV(N)$, $x\notin FV(N)$ અને
$FV\{\Subst{N}{y}{x}\}$ એમ ત્રણ જગ્યાએ મુક્ત-ચલ
સંકેતનું~$\backslash$ છૂટી ગયું છે. અહીં બધે~$\FV{}$
સંકેત વાપર્યો છે.}
\end{proof}

\begin{prob}
  \olref[lam][syn][alp]{lem:fv-one}ની સાબિતી પૂર્ણ કરો.
\end{prob}

\begin{lem}\ollabel{lem:inv}
  $P \aconvone Q$ હોય તો $Q \aconvone P$.
\end{lem}

\begin{proof}
  $P \aconvone Q$ની !!{derivation} પર આગમનથી.
  \begin{enumerate}
  \item છેલ્લો નિયમ~\olref{defn:aconvone4} હોય તો~$P$
    એ~$\lambd[x][N]$ અને~$Q$ એ
    $\lambd[y][\Subst{N}{y}{x}]$ સ્વરૂપનાં છે; અહીં
    $x \neq y$, $y \notin \FV{N}$ અને~$\Subst{N}{y}{x}$
    વ્યાખ્યાયિત છે. પ્રથમ, \olref[sub]{thm:clr} મુજબ
    $x \notin
    \FV{\Subst{N}{y}{x}}$. \olref[sub]{thm:inv} મુજબ
    $\Subst{\Subst{N}{y}{x}}{x}{y}$ વ્યાખ્યાયિત પણ છે અને
    $N$ બરાબર પણ છે. તેથી~\olref{defn:aconvone4} મુજબ
    $\lambd[y][\Subst{N}{y}{x}]
    \aconvone \lambd[x][\Subst{\Subst{N}{y}{x}}{x}{y}] =
    \lambd[x][N]$.
    \sourcecorrection{OLLAM-023}{સ્થિર અંગ્રેજી મૂળ અહીં
    સ્થાનાપન્ન પછી~$y$ મુક્ત નથી એમ કહે છે; દાખલા તરીકે
    $N=x$ હોય તો પરિણામ~$y$ જ છે. વિપરીત નામ-બદલી માટે
    જરૂરી શરત~$x\notin\FV{\Subst{N}{y}{x}}$ છે, જે
    અગાઉના પ્રમેયથી મળે છે.}
  \end{enumerate}
\end{proof}

\begin{prob}
  \olref[lam][syn][alp]{lem:inv}ની સાબિતી પૂર્ણ કરો.
\end{prob}

\begin{thm}
  $\alpha$-રૂપાંતરણ પદો પરનો સામ્ય સંબંધ છે; એટલે તે
  સ્વવાચક, સંમિત અને પરંપરિત છે.
\end{thm}

\begin{proof}
  \begin{enumerate}
  \item દરેક પદ~$M$ માટે બદ્ધ ચલોનું નામ \emph{શૂન્ય} વાર
    બદલીને~$M$માંથી~$M$ મળે છે.
  \item બદ્ધ ચલોનું નામ બદલવાની કોઈ શ્રેણીથી~$P$નું
    $Q$માં $\alpha$-રૂપાંતરણ થાય તો~$Q$માંથી એ ફેરફારો
    ઊલટા ક્રમે ઉલટાવીને (\olref{lem:inv} મુજબ)~$P$ મળે છે.
  \item એક શ્રેણીથી~$P$નું~$Q$માં અને બીજીથી~$Q$નું~$R$માં
    $\alpha$-રૂપાંતરણ થાય તો પહેલી પછી બીજી શ્રેણી
    લાગુ કરીને~$P$માંથી~$R$ મળે છે.
  \end{enumerate}
\end{proof}

હવે પછી~$M$ અને~$N$ને \emph{$\alpha$-સામ્યવાળાં},
$M
\aeq N$, ત્યારે અને માત્ર ત્યારે જ કહીએ છીએ જ્યારે~$M$નું
$N$માં $\alpha$-રૂપાંતરણ થાય. હમણાં સાબિત કર્યું તેમ,
આ ત્યારે અને માત્ર ત્યારે જ બને છે જ્યારે~$N$નું પણ
$M$માં $\alpha$-રૂપાંતરણ થાય.

\begin{thm}\ollabel{thm:fv}
  $M \aeq N$ હોય તો $\FV{M} = \FV{N}$.
\end{thm}

\begin{proof}
  \olref{lem:fv-one}માંથી સીધું મળે છે.
\end{proof}

\begin{lem}\ollabel{lem:sub:R}
  $R \aeq R'$ હોય અને~$\Subst{M}{R}{y}$ વ્યાખ્યાયિત હોય
  તો~$\Subst{M}{R'}{y}$ પણ વ્યાખ્યાયિત છે અને
  $\Subst{M}{R}{y}$ સાથે $\alpha$-સામ્યવાળું છે.
\end{lem}

\begin{proof}
  અભ્યાસપ્રશ્ન.
\end{proof}

\begin{prob}
  \olref[lam][syn][alp]{lem:sub:R} સાબિત કરો.
\end{prob}

યાદ કરો કે~\olref[sub]{sec}માં કેટલાક કિસ્સામાં સ્થાનાપન્ન
અવ્યાખ્યાયિત છે. પરંતુ પદોના બદ્ધ ચલોનું નામ
$\alpha$-રૂપાંતરણથી બદલીને સ્થાનાપન્ન વ્યાખ્યાયિત બનાવી
શકાય છે. નીચેના પ્રમેય મુજબ પરિણામમાં $\alpha$-સામ્ય
જળવાય છે:

\begin{thm}\ollabel{thm:sub}
  કોઈપણ~$M$, $R$ અને~$y$ માટે એવું~$M'$ મળે છે કે
  $M \aeq M'$ અને~$\Subst{M'}{R}{y}$ વ્યાખ્યાયિત છે.
  વધુમાં, બીજી જોડી~$M'' \aeq M$ અને~$R''$ હોય, જેમાં
  $\Subst{M''}{R''}{y}$ વ્યાખ્યાયિત તથા~$R'' \aeq R$
  હોય તો
  $\Subst{M'}{R}{y} \aeq \Subst{M''}{R''}{y}$.
\end{thm}

\begin{proof}
  $M$ના રચનાક્રમ પર આગમનથી:
  \begin{enumerate}
  \item $M$ ચલ~$z$ હોય: અભ્યાસપ્રશ્ન.
  \item $M$ એ~$\lambd[x][N]$ સ્વરૂપનું હોય. $x$ અને~$y$થી
    જુદો, $z \notin \FV{N}$ અને~$z
    \notin \FV{R}$ હોય એવો ચલ~$z$ પસંદ કરો. આગમનની
    ધારણાથી એવું~$N'$ મળે છે કે~$N' \aeq N$ અને
    $\Subst{N'}{z}{x}$ વ્યાખ્યાયિત છે. ત્યારબાદ
    \olref{defn:aconvone}\olref{defn:aconvone1} મુજબ
    $\lambd[x][N] \aeq \lambd[x][N']$ પણ છે. હવે
    \olref{defn:aconvone}\olref{defn:aconvone4} મુજબ
    $\lambd[x][N']
    \aeq \lambd[z][\Subst{N'}{z}{x}]$ છે. આ પગલું
    લઈ શકાય છે કારણ કે~$z \ne x$,
    $z \notin \FV{N'}$ અને~$\Subst{N'}{z}{x}$
    વ્યાખ્યાયિત છે. પરંતુ માત્ર~$z \neq y$ અને
    $z \notin \FV{R}$થી
    $\Subst{\lambd[z][\Subst{N'}{z}{x}]}{R}{y}$
    વ્યાખ્યાયિત છે એવું ફલિત થતું નથી: અંદરના બંધનોની
    શરતો પણ તપાસવી પડે.
    \sourcecorrection{OLLAM-024}{સ્થિર અંગ્રેજી મૂળ અહીં
    ફક્ત બહારના બાંધનાર~$z$ની અથડામણ ન હોવાથી આખું
    સ્થાનાપન્ન વ્યાખ્યાયિત છે એમ કહે છે. પરંતુ
    $N=\lambd[w][y]$, $R=w$ અને~$x,y,z,w$ જુદા હોય
    ત્યારે~$\Subst{N}{z}{x}=N$ વ્યાખ્યાયિત છે, છતાં
    $\Subst{\lambd[z][N]}{w}{y}$ અંદરના~$w$થી થતી
    અથડામણને કારણે અવ્યાખ્યાયિત છે. અગાઉથી બધાં
    જરૂરી અંદરનાં બંધનોનાં નામ તાજાં કરવાની વધારાની
    દલીલ જોઈએ; સ્થિર મૂળની આ સાબિતી અહીં અધૂરી છે.
    મૂળમાંની બે~$FV$ અભિવ્યક્તિઓને પણ~$\FV{}$ કર્યું છે.}

    વધુમાં, એ જ શરતો સંતોષતાં બીજાં~$N''$ અને~$R''$
    હોય તો સ્ત્રોત નીચેની સાંકળ આપે છે; મધ્યનાં બે
    પગલાં માત્ર આલ્ફા-સામ્ય સુધી જ માન્ય છે:
    \begin{multline*}
      \Subst{(\lambd[z][\Subst{N''}{z}{x}])}{R''}{y} \\
    \begin{aligned}
      &= \lambd[z][\Subst{\Subst{N''}{z}{x}}{R''}{y}] \\
      &\aeq \lambd[z][\Subst{\Subst{N''}{z}{x}}{R}{y}] && \text{
        \olref{lem:sub:R} મુજબ}\\
      &\aeq\lambd[z][\Subst{\Subst{N'}{z}{x}}{R}{y}]
      && \text{આગમનની ધારણાથી}\\
      &=\Subst{(\lambd[z][\Subst{N'}{z}{x}])}{R}{y}
    \end{aligned}
    \end{multline*}
    \sourcecorrection{OLLAM-025}{સ્થિર અંગ્રેજી મૂળની
    સમીકરણ-સાંકળમાં~\olref{lem:sub:R} અને આગમનનો ઉપયોગ
    કરતી જગ્યાએ સાદી સમાનતા લખી છે. સંબંધિત પદો
    સામાન્ય રીતે નામમાં જુદાં હોવાથી ત્યાં માત્ર
    $\alpha$-સામ્ય મળે છે. ઉપરનો અધૂરો વ્યાખ્યાયિતતા-
    તર્ક પૂરો કર્યા વિના આ સાંકળને સંપૂર્ણ સાબિતી
    ગણવી નહીં.}
    \item $M$ એ~$(PQ)$ સ્વરૂપનું હોય: અભ્યાસપ્રશ્ન.
  \end{enumerate}
\end{proof}

\begin{prob}
  \olref[lam][syn][alp]{thm:sub}ની સાબિતી પૂર્ણ કરો.
\end{prob}

\begin{cor}\ollabel{cor:sub}
  કોઈપણ~$M$, $R$ અને~$y$ માટે એવી જોડી~$M'$ અને~$R'$
  મળે છે કે~$M' \aeq M$, $R \aeq R'$ અને
  $\Subst{M'}{R'}{y}$ વ્યાખ્યાયિત છે. વધુમાં, બીજી
  જોડી~$M'' \aeq M$ અને~$R'' \aeq R$ હોય તથા
  $\Subst{M''}{R''}{y}$ વ્યાખ્યાયિત હોય તો
  $\Subst{M'}{R'}{y} \aeq
  \Subst{M''}{R''}{y}$.
  \sourcecorrection{OLLAM-026}{સ્થિર અંગ્રેજી મૂળના
  ઉપસિદ્ધાંતમાં બીજી જોડી માટે~$R''\aeq R$ની શરત
  ગાયબ છે અને તેની જગ્યાએ પહેલી જોડીનું સ્થાનાપન્ન
  ફરી વ્યાખ્યાયિત હોવાનું લખાયું છે. અહીં બીજા
  સ્થાનાપન્નની વ્યાખ્યાયિતતા તથા બંને પ્રતિસ્થાપકોનું
  $\alpha$-સામ્ય સ્પષ્ટ કર્યું છે.}
\end{cor}

\begin{proof}
  \olref{thm:sub}માંથી સીધું મળે છે.
\end{proof}

\end{document}
